\podnaslov{Neodvisnost slučajnih spremenljivk}

\begin{definicija}
  Slučajni spremenljivki $X$ in $Y$ sta \pojem{neodvisni}, če za vsaki Borelovi
  $A$ in $B$ velja $P(X \in A, Y \in B) = P(X \in A) P(Y \in B)$.
\end{definicija}

\begin{definicija}
  Slučajne spremenljivke $X_1, \ldots, X_n$ so \pojem{neodvisne}, če za vsak
  nabor Borelovih množic $A_1, \ldots, A_n$ velja
  \[
	P(\forall i \velja X_i \in A_i) = \prod_j P(X_j \in A_j).
  \]
\end{definicija}

\begin{definicija}
  Slučajne spremenljivke $\{X_i\}_{i \in I}$ so \pojem{neodvisne}, če so
  neodvisne vse končne poddružine.
\end{definicija}

\vprasanje{Definiraj neodvisnost slučajnih spremenljivk.}

